\documentclass[Orbiter User Manual.tex]{subfiles}
\begin{document}

\section{Command line interface}
When starting Orbiter (orbiter.exe) or the Orbiter graphics server
(orbiter\_ng.exe) from the command line, optional parameters can be provided to
modify the behaviour of the simulation.

\begin{lstlisting}[language=OSFS]
orbiter.exe [options]
\end{lstlisting}

Below is a list of the supported options. Options are separated with space. Most
options can be specified with a long notation (with prefix -{}-) or short
notation (with prefix -). If the option requires an argument, it must be
separated from the option key with = (for the long version) or with a space (for
the short version). If an argument contains spaces, it should be enclosed in
double quotes (").

%\begin{table}[H]
	%\centering
	\begin{longtable}{ p{0.2\textwidth} p{0.12\textwidth}
		>{\RaggedRight}p{0.6\textwidth} }
	\hline
	\textbf{Long} & \textbf{Short} & \textbf{Effect}\\
	\hline
	-{}-help & -h & Display usage and list of options, then exit.\\
	\hline
	-{}-cfgcinonly & -I & Read configuration from stdin instead of Orbiter.cfg,
		e.g.
\begin{lstlisting}[language=OSFS,aboveskip=2pt, belowskip=2pt]
Orbiter.exe -I < Local.cfg.
\end{lstlisting}
		Implies -o. (See notes below)\\
	\hline
	-{}-cfgcinadd & -i & Read configuration from stdin after reading
		Orbiter.cfg, e.g.
\begin{lstlisting}[language=OSFS,aboveskip=2pt, belowskip=2pt]
echo EnableClouds = FALSE | Orbiter.exe -i.
\end{lstlisting}\\
	\hline
	-{}-cfgreadonly & -o & Do not update Orbiter.cfg on exit. Forget all
		configuration changes since starting Orbiter.\\
	\hline
	-{}-cfgcoutonly & -O & Write configuration to stdout instead of Orbiter.cfg.
		(See notes below.)\\
	\hline
	-{}-cfgname=\textit{<name>} & -n \textit{<name>} & Use \textit{<name>}
		instead of Orbiter.cfg for configuration I/O.\\
	\hline
	-{}-fastexit & -x & Terminate Orbiter at the end of the session. Don't
		return to the Launchpad dialog.\\
	\hline
	-{}-paused & -P & Force the simulation to start paused, bypassing Launchpad
		setting.\\
	\hline
	-{}-unpaused & -R & Force the simulation to start unpaused, bypassing
		Launchpad setting.\\
	\hline
	-{}-scenario=\textit{<scn>} & -S \textit{<scn>} & Directly launch scenario
		\textit{<scn>}. Bypass the Launchpad dialog. \textit{<scn>} should
		contain the path to the scenario file relative to .\textbackslash
		Scenarios. The file extension (.scn) should be omitted.\\
	\hline
	-{}-scenariox=\textit{<scn>} & -s \textit{<scn>} & Shortcut for -x -S
		\textit{<scn>}. Launch scenario directly and terminate on session end.\\
	\hline
	-{}-fastexit & -x & Exit Orbiter after simulation session.\\
	\hline
	-{}-openvideotab & -v & Open Launchpad dialog on video tab instead of
		scenarios.\\
	\hline
	-{}-nographics & -G & Launch Orbiter in console mode without loading
		graphics clients.\\
	\hline
	-{}-keeplog & -l & Append log output to Orbiter.log instead of overwriting
		the previous log data.\\
	\hline
	-{}-fixedstep=\textit{<s>} & -f \textit{<s>} & Set a fixed simulation time
		step of \textit{<s>} seconds for each frame, independent of how long the
		frame takes to render in real time.\\
	\hline
	-{}-maxsystime=\textit{<t>} & -T \textit{<t>} & Terminate the simulation
		session after \textit{<t>} seconds.\\
	\hline
	-{}-maxsimtime=\textit{<t>} & -t \textit{<t>} & Terminate the simulation
		session at simulation time \textit{<t>}. Implies -R unless explicitly
		overridden with -P.\\
	\hline
	-{}-maxframes=\textit{<f>} & & Terminate the simulation session after
		\textit{<f>} time frames. Implies -R unless explicitly overridden with
		-P.\\
	\hline
	-{}-plugin=\textit{<pg>} & -p \textit{<pg>} & Enforce loading of plugin
		\textit{<pg>}. Any path provided must be relative to .\textbackslash
		Modules\textbackslash Plugin. The extension (.dll) should be omitted.
		Multiple -{}-plugin options can be provided. Any plug-ins requested on
		the command line cannot be unloaded interactively.\\
	\hline
	\end{longtable}
%\end{table}

\textbf{Note on input and output redirection:}
When using "<" and ">" for redirecting input and output to and from a configuration file, please be aware that because of the way the operating system
processes the command line, you cannot use the same file for input and output, i.e. you cannot use
\begin{lstlisting}[language=OSFS]
Orbiter.exe -I -O < Local.cfg > Local.cfg
\end{lstlisting}
because this will truncate the configuration file before Orbiter has the chance to read it. Instead use
\begin{lstlisting}[language=OSFS]
Orbiter.exe -I -O < Local.cfg > Local.tmp && move /y Local.tmp Local.cfg
\end{lstlisting}
or just use the -{}-cfgname (-n) option:
\begin{lstlisting}[language=OSFS]
Orbiter.exe -n Local.cfg
\end{lstlisting}

\textbf{Interactive parameter input:}
If you redirect configuration input from stdin (with -i or -I), but don't redirect from a file or pipe on the command line, Orbiter will wait for
manual input from the console. You can then type in parameters manually (using the same format as used in the Orbiter.cfg file). Terminate
your input with an end-of-file signal (\keystroke{\ctrl,Z} for Windows).

\end{document}
